% Source: Wald GR, physical PDF pp. 49--54
%         (printed pp. 36--41).
% Coverage: Section 3.2, Curvature; equations (3.2.1)--(3.2.32).
% Figures/tables/footnotes: Figure 3.3; no tables or footnotes.
% Status: source-reviewed and compile-clean.
% Uncertainties: none; a printed-source typo is corrected below.

\subsection{Curvature}\label{sec:3.2}

In the previous section we showed that given a derivative operator, there exists a
notion of how to parallel transport a vector from $p$ to $q$ along a curve $C$. However,
the vector in $T_qM$ which we get by this parallel transport procedure starting from a
vector in $T_pM$ will, in general, depend on the choice of curve connecting them. As
already indicated in the discussion at the beginning of this chapter, we can use the
path dependence of parallel transport to define an intrinsic notion of curvature. In this
section we carry out this program by first defining the Riemann curvature tensor in
terms of the failure of successive operations of differentiation to commute when
applied to a covector field. Then we show that this tensor is directly related to the
path-dependent nature of parallel transport; specifically, the failure of a vector to
return to its original value when parallel transported around a small closed loop is
governed by the Riemann tensor. In the next section we will show that the Riemann
tensor also fully describes the other characterization of curvature mentioned above:
the failure of initially parallel geodesics to remain parallel.

Let $\nabla_a$ be a derivative operator. Let $\omega_a$ be a covector field and let $f$ be a smooth
function. We calculate the action of two derivative operators applied to $f\omega_a$,
\begin{equation}
  \begin{aligned}
    \nabla_a\nabla_b(f\omega_c)
    &=\nabla_a(\omega_c\nabla_bf+f\nabla_b\omega_c)\\
    &=(\nabla_a\nabla_bf)\omega_c+\nabla_bf\nabla_a\omega_c
      +\nabla_af\nabla_b\omega_c+f\nabla_a\nabla_b\omega_c.
  \end{aligned}
  \tag{3.2.1}\label{eq:3.2.1}
\end{equation}
If we subtract from this the tensor $\nabla_b\nabla_a(f\omega_c)$, the first three terms of the right-hand
side of Eq.~\eqref{eq:3.2.1} will cancel the corresponding terms of the expression for
$\nabla_b\nabla_a(f\omega_c)$ and we obtain the simple result,
\begin{equation}
  (\nabla_a\nabla_b-\nabla_b\nabla_a)(f\omega_c)
  =f(\nabla_a\nabla_b-\nabla_b\nabla_a)\omega_c.
  \tag{3.2.2}\label{eq:3.2.2}
\end{equation}
By exactly the same reasoning as given above in the discussion of derivative operators (see Eq.~\eqref{eq:3.1.3}), it follows that the tensor
$(\nabla_a\nabla_b-\nabla_b\nabla_a)\omega_c$ at point $p$ depends
only on the value of $\omega_c$ at $p$. Consequently,
$(\nabla_a\nabla_b-\nabla_b\nabla_a)$ defines a linear map from
covectors at $p$ to type $(0,3)$ tensors at $p$; i.e., its action is that of a tensor of type
$(1,3)$. Thus, we have shown that there exists a tensor field $R_{abc}{}^d$ such that for all
covector fields $\omega_c$, we have
\begin{equation}
  \nabla_a\nabla_b\omega_c-\nabla_b\nabla_a\omega_c=R_{abc}{}^d\omega_d.
  \tag{3.2.3}\label{eq:3.2.3}
\end{equation}
$R_{abc}{}^d$ is called the \emph{Riemann curvature tensor}.

We first show that $R_{abc}{}^d$ is directly related to the failure of a vector to return to its
initial value when parallel transported around a small closed curve. We can conveniently construct a small closed loop at $p\in M$ by choosing a two-dimensional
surface $S$ through $p$ and choosing coordinates $t$ and $s$ in the surface [with the
coordinates of $p$ chosen, for simplicity, to be $(0,0)$]. Consider the loop formed by
moving $\Delta t$ along the $s=0$ curve, followed by moving $\Delta s$ along the $t=\Delta t$ curve,
and then moving back by $\Delta t$ and $\Delta s$ as illustrated in Fig.~\ref{fig:3.3}. Let $v^a$ be a vector
at $p$ (not necessarily tangent to $S$) and let us parallel transport $v^a$ around this closed

% Source: Wald GR, physical PDF p. 50 (printed p. 37).
% Coverage: Figure 3.3.
% Figures/tables/footnotes: Figure only; no tables or footnotes.
% Status: source-reviewed and compile-clean.
% Uncertainties: none.
\begingroup
\renewcommand{\thefigure}{3.3}
\begin{figure}[H]
  \centering
  \begin{tikzpicture}[scale=.92,>=Stealth]
    \draw[thick] (-3.0,-0.90) .. controls (-2.1,0.55) and (-0.6,1.35) .. (2.60,1.00)
      -- (1.95,-1.05) .. controls (-0.3,-0.62) and (-1.9,-1.35) .. (-3.0,-0.90);
    \node at (-2.15,-0.50) {$S$};
    \coordinate (P) at (-.95,-.34);
    \coordinate (A) at (.22,-.10);
    \coordinate (B) at (.62,.72);
    \coordinate (C) at (-.55,.49);
    \fill (P) circle (1.6pt);
    \draw[->,thick] (P) -- (A);
    \draw[->,thick] (A) -- (B);
    \draw[->,thick] (B) -- (C);
    \draw[->,thick] (C) -- (P);
    \node[below left=2pt] at (P) {$p$};
    \node[below right=3pt] at (A) {$(\Delta t,0)$};
    \node[above right=2pt] at (B) {$(\Delta t,\Delta s)$};
    \node[above left=2pt] at (C) {$(0,\Delta s)$};
    % The transported vector at successive stages: over a small loop it barely
    % changes direction, the residual rotation being the Riemann tensor at p.
    \foreach \x/\y in {-.95/-.34,.22/-.10,.62/.72,-.55/.49} {
      \draw[->] (\x,\y) -- ++(.20,.28);
    }
    \node[left=3pt] at (-.88,-.03) {$v^a$};
  \end{tikzpicture}
  \caption{The parallel transport of a vector $v^a$ around a small closed loop. As
  derived in the text, to second order in $\Delta t$ and $\Delta s$, the change in $v^a$ is governed by
  the Riemann tensor at $p$.}
  \label{fig:3.3}
\end{figure}
\endgroup


loop. It is easiest to compute the change in $v^a$ when we return to $p$ by letting $\omega_a$ be
an arbitrary covector field and finding the change in the scalar $v^a\omega_a$ as we traverse
the loop. For small $\Delta t$ the change, $\delta_1$, in $v^a\omega_a$ on the first leg of the path is
\begin{equation}
  \delta_1=\left.\Delta t\frac{\partial}{\partial t}(v^a\omega_a)\right|_{(\Delta t/2,0)},
  \tag{3.2.4}\label{eq:3.2.4}
\end{equation}
where, by evaluating the derivative at the midpoint, we have ensured that this
expression is accurate to second order in the displacement $\Delta t$. We may rewrite $\delta_1$ as
\begin{equation}
  \begin{aligned}
    \delta_1
    &=\left.\Delta t\,T^b\nabla_b(v^a\omega_a)\right|_{(\Delta t/2,0)}\\
    &=\left.\Delta t\,v^aT^b\nabla_b\omega_a\right|_{(\Delta t/2,0)},
  \end{aligned}
  \tag{3.2.5}\label{eq:3.2.5}
\end{equation}
where $T^b$ is the tangent to the curves of constant $s$ and $T^b\nabla_bv^a=0$ by
Eq.~\eqref{eq:3.1.16}. Similar expressions hold for the changes $\delta_2$, $\delta_3$, and
$\delta_4$ on the other parts of the path. The two \lq\lq $\Delta t$ variations,\rq\rq\ $\delta_1$ and
$\delta_3$, combine to yield
\begin{equation}
  \delta_1+\delta_3
  =\Delta t\left\{
  \left.v^aT^b\nabla_b\omega_a\right|_{(\Delta t/2,0)}
  -\left.v^aT^b\nabla_b\omega_a\right|_{(\Delta t/2,\Delta s)}
  \right\},
  \tag{3.2.6}\label{eq:3.2.6}
\end{equation}
and $\delta_2$ and $\delta_4$ combine similarly. Since the term in brackets vanishes as
$\Delta s\to0$, this shows that to first order in $\Delta t$ and $\Delta s$, the total change in $v^a\omega_a$
(and thus the total change in $v^a$) vanishes; i.e., parallel transport is path-independent
to \emph{first order} in $\Delta t$ and $\Delta s$. To calculate the second order change in $v^a\omega_a$, we need to evaluate the
term in brackets in Eq.~\eqref{eq:3.2.6} to first order. We do this by the following
procedure: We consider the curve $t=\Delta t/2$ and imagine parallel transporting $v^a$ and
$T^b\nabla_b\omega_a$ along this curve from $(\Delta t/2,0)$ to $(\Delta t/2,\Delta s)$. Now to first order in
$\Delta s$, $v^a$ at $(\Delta t/2,\Delta s)$ equals the parallel transport of $v^a$ at $(\Delta t/2,0)$ along this curve since,
as remarked above, parallel transport is path-independent to first order. On the other
hand, to first order, the term $T^b\nabla_b\omega_a$ at $(\Delta t/2,\Delta s)$ will differ from the parallel
transport of that quantity from $(\Delta t/2,0)$ by the amount
$\Delta sS^c\nabla_c(T^b\nabla_b\omega_a)$, where $S^c$ is the tangent to the curves of constant $t$. Hence, the term in
brackets is just this quantity contracted with $v^a$. Thus, to second order in $\Delta t$, $\Delta s$, we find
\begin{equation}
  \delta_1+\delta_3=-\Delta t\Delta s\,v^aS^c\nabla_c(T^b\nabla_b\omega_a),
  \tag{3.2.7}\label{eq:3.2.7}
\end{equation}
where, to this accuracy, we may evaluate all tensors at $p$. Adding the similar
contribution for $\delta_2$ and $\delta_4$, we find the total change in $v^a\omega_a$ is
\begin{equation}
  \begin{aligned}
    \delta(v^a\omega_a)
    &={}\Delta t\Delta s\,v^a\bigl\{T^c\nabla_c(S^b\nabla_b\omega_a)
      -S^c\nabla_c(T^b\nabla_b\omega_a)\bigr\}\\
    &={}\Delta t\Delta s\,v^aT^cS^b(\nabla_c\nabla_b-\nabla_b\nabla_c)\omega_a\\
    &={}\Delta t\Delta s\,v^aT^cS^bR_{cba}{}^d\omega_d,
  \end{aligned}
  \tag{3.2.8}\label{eq:3.2.8}
\end{equation}
where in the second line we used the fact that the coordinate vector fields $T^a$ and $S^b$
commute (see the end of section 2.2 and Eq.~\eqref{eq:3.1.2}) and we used the definition of
the Riemann tensor Eq.~\eqref{eq:3.2.3} in the last step. But Eq.~\eqref{eq:3.2.8} can hold
for all $\omega_a$ if and only if the total change in $v^a$ (accurate to second order in $\Delta t$ and $\Delta s$)
is
\begin{equation}
  \delta v^a=\Delta t\Delta s\,v^dT^cS^bR_{cbd}{}^a.
  \tag{3.2.9}\label{eq:3.2.9}
\end{equation}
This is the desired result. It shows that the Riemann tensor indeed directly measures
the path dependence of parallel transport.

Using Eq.~\eqref{eq:3.2.3} we may---by a procedure completely analogous to the
derivation of Eq.~\eqref{eq:3.1.14}---express the action of the commutator of derivative
operators on arbitrary tensor fields in terms of the Riemann tensor. To find the
expression for a vector field $t^a$, we let $\omega_a$ be a covector field. Then using property
5 of derivative operators together with the Leibniz rule and Eq.~\eqref{eq:3.2.3}, we find
\begin{equation}
  \begin{aligned}
    0
    &=(\nabla_a\nabla_b-\nabla_b\nabla_a)(t^c\omega_c)\\
    &=\nabla_a(\omega_c\nabla_bt^c+t^c\nabla_b\omega_c)
      -\nabla_b(\omega_c\nabla_at^c+t^c\nabla_a\omega_c)\\
    &=\omega_c(\nabla_a\nabla_b-\nabla_b\nabla_a)t^c
      +t^c(\nabla_a\nabla_b-\nabla_b\nabla_a)\omega_c\\
    &=\omega_c(\nabla_a\nabla_b-\nabla_b\nabla_a)t^c+t^c\omega_dR_{abc}{}^d.
  \end{aligned}
  \tag{3.2.10}\label{eq:3.2.10}
\end{equation}
Thus, we obtain
\begin{equation}
  (\nabla_a\nabla_b-\nabla_b\nabla_a)t^c=-R_{abd}{}^ct^d.
  \tag{3.2.11}\label{eq:3.2.11}
\end{equation}
By induction, for an arbitrary tensor field $T^{c_1\cdots c_k}{}_{d_1\cdots d_l}$ we find
\begin{equation}
  \begin{aligned}
    (\nabla_a\nabla_b-\nabla_b\nabla_a)T^{c_1\cdots c_k}{}_{d_1\cdots d_l}
    ={}&-\sum_{i=1}^{k}R_{abe}{}^{c_i}
    T^{c_1\cdots e\cdots c_k}{}_{d_1\cdots d_l}\\
    &+\sum_{j=1}^{l}R_{abd_j}{}^e
    T^{c_1\cdots c_k}{}_{d_1\cdots e\cdots d_l}.
  \end{aligned}
  \tag{3.2.12}\label{eq:3.2.12}
\end{equation}

Next we establish four key properties of the Riemann tensor:
\begin{enumerate}[label=\arabic*.,leftmargin=*]
\item
\begin{equation}
  R_{abc}{}^d=-R_{bac}{}^d.
  \tag{3.2.13}\label{eq:3.2.13}
\end{equation}
\item
\begin{equation}
  R_{[abc]}{}^d=0.
  \tag{3.2.14}\label{eq:3.2.14}
\end{equation}
\item For the derivative operator $\nabla_a$ naturally associated with the metric, $\nabla_ag_{bc}=0$, we
have
\begin{equation}
  R_{abcd}=-R_{abdc}.
  \tag{3.2.15}\label{eq:3.2.15}
\end{equation}
\item The Bianchi identity holds:
\begin{equation}
  \nabla_{[a}R_{bc]d}{}^e=0.
  \tag{3.2.16}\label{eq:3.2.16}
\end{equation}
\end{enumerate}

Property (1) follows trivially from the definition of $R_{abc}{}^d$, Eq.~\eqref{eq:3.2.3}. To
prove property (2), we note that for an arbitrary covector field, $\omega_a$, and any
derivative operator $\nabla_a$, we have
\begin{equation}
  \nabla_{[a}\nabla_b\omega_{c]}=0.
  \tag{3.2.17}\label{eq:3.2.17}
\end{equation}
This equation can be proven from Eq.~\eqref{eq:3.1.14}, substituting an ordinary derivative $\partial_a$ for $\widetilde\nabla_a$, and using the commutativity of ordinary derivatives and the symmetry
of $C^c{}_{ab}=\Gamma^c{}_{ab}$, Eq.~\eqref{eq:3.1.9}. (In differential forms notation, it is the statement
that $d^2\omega=0$ [see appendix B].) Thus, we have
\begin{equation}
  0=2\nabla_{[a}\nabla_b\omega_{c]}
  =\nabla_{[a}\nabla_b\omega_{c]}-\nabla_{[b}\nabla_a\omega_{c]}
  =R_{[abc]}{}^d\omega_d
  \tag{3.2.18}\label{eq:3.2.18}
\end{equation}
for all $\omega_d$, which proves property (2).

Property (3) follows from Eq.~\eqref{eq:3.2.12} applied to the metric $g_{ab}$. We find
\begin{equation}
  0=(\nabla_a\nabla_b-\nabla_b\nabla_a)g_{cd}
  =R_{abc}{}^eg_{ed}+R_{abd}{}^eg_{ce}
  =R_{abcd}+R_{abdc},
  \tag{3.2.19}\label{eq:3.2.19}
\end{equation}
which yields property (3). It follows from properties 1, 2, and 3 that the Riemann
tensor also satisfies the following useful symmetry property (see problem 3):
\begin{equation}
  R_{abcd}=R_{cdab}.
  \tag{3.2.20}\label{eq:3.2.20}
\end{equation}

Finally, to prove the Bianchi identity, property (4), we apply the commutator of
derivative operators to the derivative of a covector field. We obtain, using Eq.~\eqref{eq:3.2.12},
\begin{equation}
  (\nabla_a\nabla_b-\nabla_b\nabla_a)\nabla_c\omega_d
  =R_{abc}{}^e\nabla_e\omega_d+R_{abd}{}^f\nabla_c\omega_f.
  \tag{3.2.21}\label{eq:3.2.21}
\end{equation}
On the other hand, we have
\begin{equation}
  \begin{aligned}
    \nabla_a(\nabla_b\nabla_c\omega_d-\nabla_c\nabla_b\omega_d)
    &=\nabla_a(R_{bcd}{}^e\omega_e)\\
    &=\omega_e\nabla_aR_{bcd}{}^e+R_{bcd}{}^e\nabla_a\omega_e.
  \end{aligned}
  \tag{3.2.22}\label{eq:3.2.22}
\end{equation}
If we antisymmetrize over $a$, $b$, and $c$ in Eqs.~\eqref{eq:3.2.21} and~\eqref{eq:3.2.22}, the
left-hand sides become equal. Equality of the right-hand sides yields
\begin{equation}
  R_{[abc]}{}^e\nabla_e\omega_d
  +R_{[ab|d|}{}^e\nabla_{c]}\omega_e
  =\omega_e\nabla_{[a}R_{bc]d}{}^e
  +R_{[bc|d|}{}^e\nabla_{a]}\omega_e,
  \tag{3.2.23}\label{eq:3.2.23}
\end{equation}
where the vertical bars indicate that we do not antisymmetrize over $d$. The first term
on the left-hand side vanishes by Eq.~\eqref{eq:3.2.14} while the second terms on both
sides cancel each other. Thus, we obtain, for all $\omega_e$,
\begin{equation}
  \omega_e\nabla_{[a}R_{bc]d}{}^e=0,
  \tag{3.2.24}\label{eq:3.2.24}
\end{equation}
which yields property (4).

It is useful to decompose the Riemann tensor into a \lq\lq trace part\rq\rq\ and a \lq\lq trace-free
part.\rq\rq\ % SOURCE ERRATUM: printed ``trace tree part.''
By the antisymmetry properties (1) and (3), the trace of the Riemann tensor
over its first two or last two indices vanishes. However, its trace over the second and
fourth (or equivalently, the first and third) indices defines the \emph{Ricci tensor}, $R_{ac}$,
\begin{equation}
  R_{ac}=R_{abc}{}^b.
  \tag{3.2.25}\label{eq:3.2.25}
\end{equation}
By Eq.~\eqref{eq:3.2.20}, $R_{ab}$ satisfies the symmetry property
\begin{equation}
  R_{ac}=R_{ca}.
  \tag{3.2.26}\label{eq:3.2.26}
\end{equation}
The \emph{scalar curvature}, $R$, is defined as the trace of the Ricci tensor:
\begin{equation}
  R=R_a{}^a.
  \tag{3.2.27}\label{eq:3.2.27}
\end{equation}
The \lq\lq trace free part\rq\rq\ is called the \emph{Weyl tensor}, $C_{abcd}$, and is defined for manifolds of
dimension $n\geq3$ by the equation
\begin{equation}
  R_{abcd}=C_{abcd}+\frac{2}{n-2}
  \bigl(g_{a[c}R_{d]b}-g_{b[c}R_{d]a}\bigr)
  -\frac{2}{(n-1)(n-2)}R\,g_{a[c}g_{d]b}.
  \tag{3.2.28}\label{eq:3.2.28}
\end{equation}
The Weyl tensor satisfies the symmetry properties (1), (2), and (3) of the Riemann
tensor as well as being trace free on all its indices. It also behaves in a very simple
manner under conformal transformations of the metric (see appendix D) and for this
reason is sometimes called the \emph{conformal tensor}.

Contraction of the Bianchi identity~\eqref{eq:3.2.16} leads to an important equation
satisfied by $R_{ab}$. We find
\begin{equation}
  \nabla_aR_{bcd}{}^a+\nabla_bR_{cd}-\nabla_cR_{bd}=0.
  \tag{3.2.29}\label{eq:3.2.29}
\end{equation}
Raising the index $d$ with the metric and contracting over $b$ and $d$, we obtain
\begin{equation}
  \nabla_aR_c{}^a+\nabla_bR_c{}^b-\nabla_cR=0
  \tag{3.2.30}\label{eq:3.2.30}
\end{equation}
or
\begin{equation}
  \nabla^aG_{ab}=0,
  \tag{3.2.31}\label{eq:3.2.31}
\end{equation}
where
\begin{equation}
  G_{ab}=R_{ab}-\frac{1}{2}Rg_{ab}.
  \tag{3.2.32}\label{eq:3.2.32}
\end{equation}
The tensor $G_{ab}$ is called the \emph{Einstein tensor}. It appears in Einstein's equation, and the
twice contracted Bianchi identity, Eq.~\eqref{eq:3.2.31}, plays an important role in
ensuring consistency of Einstein's equation (see chapters 4 and 10).
